\section{总结与延伸}

本节把整期压缩成五个结论。第一，Agentic LLM 的关键是从“缸中之脑”走向能使用工具的数字行动者。第二，test-time scaling 同时包括长思考和 Agent 工具反馈。第三，K2 把基础模型、token efficiency、Muon/Rephrase、Agentic 泛化和开源生态放在同一套设计中。第四，数据墙之后，FLOPs、RL、反馈和架构效率会重新平衡，Linear Attention 等低成本路线必须警惕能力 bias。第五，模型公司不仅要解技术题，也要解开源、商业模式、组织反馈和创始人心态题。

\begin{knowledgebox}{关键判断清单}
\begin{tabularx}{\textwidth}{p{0.24\textwidth}p{0.34\textwidth}X}
\toprule
判断 & 为什么成立 & 仍需观察 \\
\midrule
Agentic 是下一个主线 & 模型已经能写代码、用工具、长程执行 & 泛化和可靠性是否能持续提升。 \\
Token efficiency 变重要 & 高质量数据边际变稀缺 & 数据改写和优化器是否稳定。 \\
开源有战略意义 & 生态验证和开发者扩散很强 & 商业捕获和维护成本。 \\
架构效率不能牺牲智商 & 长上下文成本必须降 & Linear/Sparse/Hybrid 的能力上限。 \\
商业边界会重画 & API、Agent 产品、生态互相牵制 & 谁掌握用户入口和反馈数据。 \\
\bottomrule
\end{tabularx}
\end{knowledgebox}

\begin{importantbox}{把 EP113 放进张小珺 AI 队列}
EP119 讲 Attention 架构考古，EP115 讲 Agent 下半场理论，EP116 讲企业级 Agentic Model；EP113 则给出模型公司视角：K2 如何把基础模型、Agentic 能力和开源生态结合，并在数据墙后寻找新的 scaling 路线。
\end{importantbox}

\subsection{关键 takeaways}

\begin{enumerate}
\item K2 的价值在于基础模型能力、编程/工具使用、Agentic 泛化和开源生态的组合。
\item Test-time scaling 是长思考与 Agent 工具使用的共同抽象。
\item Token efficiency 是数据墙时代的核心训练指标之一。
\item Linear Attention 等长上下文架构要同时看效率和“智商”风险。
\item 开源、API、Agent 产品和基座模型边界，是模型公司商业化的关键问题。
\end{enumerate}

\subsection{与前后几集的关系}

\begin{tabularx}{\textwidth}{p{0.18\textwidth}p{0.34\textwidth}X}
\toprule
节目 & 主题 & 与 EP113 的连接 \\
\midrule
EP119 & Kimi Linear、DeepSeek、MiniMax 注意力架构 & 解释 K2 背后的 Attention/长上下文架构取舍。 \\
EP115 & Agent 下半场、reward、interface & 给出 Agentic LLM 的理论框架。 \\
EP116 & 企业级 Agentic Model & 展示 Agentic 能力在 ToB 私有数据场景的落地版本。 \\
EP118 & VLA、Agent OS、物理世界入口 & 把“模型长出手”从数字世界推向物理世界。 \\
\bottomrule
\end{tabularx}

\subsection{开放问题}

前面给出结论，本节保留真正没有定论的部分。它们不是结尾装饰，而是 K2 之后模型公司和 Agent 产品公司都要继续回答的问题。

\begin{enumerate}
\item 数据墙之后，FLOPs scaling、RL feedback 和架构效率会怎样重新分配预算？
\item Linear Attention、Sparse Attention、Hybrid Attention 哪条路线能在能力和成本之间胜出？
\item 开源模型如何在生态扩散和商业捕获之间找到平衡？
\item Agentic LLM 的泛化应该如何评测，特别是在未知工具和未知任务上？
\item 基座模型公司会吞掉 Agent 应用，还是 Agent 应用会反过来定义基座模型需求？
\end{enumerate}

\begin{warningbox}{开放问题的价值}
这些问题不是附录，而是模型公司真实决策的核心。K2 的意义也正在于它把这些问题同时暴露出来：模型、数据、工具、生态、商业和组织没有一个可以单独解决。
\end{warningbox}

\begin{knowledgebox}{术语对照附表}
\begin{tabularx}{\textwidth}{p{0.22\textwidth}p{0.34\textwidth}X}
\toprule
术语 & 可记忆解释 & 在本期中的位置 \\
\midrule
缸中之脑 & 只有认知、没有行动接口的模型 & 导读和 Agentic LLM 主线。 \\
万能构造器 & 能利用工具改变世界的行动者 & Agent 能力和人类类比。 \\
数据墙 & 高质量训练数据边际变稀缺 & Scaling 与数据飞轮讨论。 \\
智商风险 & 降成本架构可能损害模型能力 & Linear Attention 取舍。 \\
RL 管理 & 用目标和反馈塑造团队探索 & 创始人组织隐喻。 \\
\bottomrule
\end{tabularx}
\end{knowledgebox}

\subsection{拓展阅读}

\begin{itemize}
\item 对 Attention 架构取舍感兴趣，可对照 EP119 的 Kimi Linear / MiniMax M2 / DeepSeek 注意力综述。
\item 对 Agent 理论感兴趣，可对照 EP115 姚顺雨访谈，理解 reward、multi-agent 和 interface。
\item 对企业落地感兴趣，可对照 EP116 吴明辉访谈，观察 Agentic Model 在 ToB 私有数据场景中的形态。
\end{itemize}

\end{document}
